\chaptercustom{1}{CÁC QUY ĐỊNH CHUNG}
Chương này tóm tắt các quy định trình bày đồ án tốt nghiệp. Các chương từ 2 đến 5 hướng dẫn cách dùng mẫu LaTeX này: mỗi tính năng được giải thích và dùng ngay trên chính trang đó để bạn xem kết quả. Khi viết đồ án, hãy xóa nội dung mẫu và thay bằng nội dung của bạn.

\subsection{Bố cục đồ án}
Phần mở đầu giới thiệu vấn đề mà đồ án cần giải quyết, mô tả các phương pháp hiện có để giải quyết vấn đề, trình bày mục đích của đồ án song song với việc giới hạn phạm vi của vấn đề mà đồ án tập trung giải quyết. Phần này cũng giới thiệu tóm tắt cấu trúc đồ án và nội dung tương ứng của các phần sẽ lần lượt được trình bày ở các chương tiếp theo.

Nội dung chính của một đồ án tốt nghiệp bao gồm:
\begin{itemize}
    \item Phần mở đầu giới thiệu đề tài.
    \item Một chương giới thiệu cơ sở lý thuyết.
    \item Một hoặc nhiều chương trình bày các vấn đề về tính toán và thiết kế.
    \item Một chương mô tả các thí nghiệm và kết quả thu được.
\end{itemize}

Mỗi chương bắt đầu bằng đoạn giới thiệu các phần chính được trình bày trong chương đó, dài khoảng 5--10 dòng, và kết thúc bằng một đoạn tóm tắt các kết luận chính của chương. Chú ý phân bổ chiều dài của mỗi chương cho cân đối và hợp lý.

\subsection{Nộp đồ án}
Sinh viên (hoặc nhóm sinh viên tối đa 3 thành viên làm chung một đề tài) phải nộp 2 quyển đồ án tốt nghiệp tại văn phòng bộ môn của giảng viên hướng dẫn trước ngày bảo vệ ít nhất 1 tuần. Một quyển đồ án cần có các đặc điểm sau:
\begin{itemize}
    \item Được \textbf{in 2 mặt} nhằm tiết kiệm không gian lưu trữ.
    \item Đóng bìa mềm, bên ngoài là bóng kính.
    \item Số trang 50--150 trang, không kể phần phụ lục.
    \item Phải có chữ ký của sinh viên sau lời cam đoan và của giảng viên hướng dẫn.
\end{itemize}

\subsection{Phụ lục}
Phụ lục (nếu có) chứa thông tin có liên quan đến đồ án nhưng nếu để trong phần chính sẽ gây rườm rà. Thông thường các chi tiết để trong phần phụ lục là kết quả thô (chưa qua xử lý), mã nguồn phần mềm, thông số chi tiết của linh kiện hoặc hình ảnh minh họa thêm.

\subsection{Tài liệu tham khảo}
\subsubsection{Cách liệt kê}
Áp dụng cách liệt kê theo quy định của IEEE. Theo đó tài liệu tham khảo được đánh số thứ tự trong ngoặc vuông. Thứ tự liệt kê là thứ tự xuất hiện của tài liệu tham khảo được trích dẫn trong đồ án. Tài liệu tham khảo đã liệt kê bắt buộc phải được trích dẫn trong phần nội dung của đồ án. Tài liệu tham khảo cần có nguồn gốc rõ ràng và phải từ nguồn đáng tin cậy. Cần hạn chế trích dẫn tài liệu tham khảo từ các website, từ Wikipedia. Mẫu này tự làm đúng các quy định trên (xem Mục~\ref{sec:trichdan}).

\subsubsection{Các loại tài liệu tham khảo}
Các nguồn tài liệu tham khảo chính là sách, bài báo trong các tạp chí, bài báo trong các hội nghị khoa học và các tài liệu tham khảo khác trên internet.

\subsection{Đánh số hình vẽ, bảng biểu, phương trình}
Hình vẽ, bảng biểu và phương trình được đánh số theo số của chương, ví dụ Hình 3.1, Bảng 3.2, PT 3.3. Chú thích của hình vẽ đặt ngay dưới hình, chú thích của bảng đặt ngay trên bảng. Mọi hình vẽ và bảng biểu đều phải được nhắc tới, phân tích và bình luận trong phần nội dung.

Các định nghĩa, định lý, hệ quả, bổ đề được đánh số theo số của chương và dùng chung một dãy số. Ví dụ trong chương 3, chúng được đánh số theo thứ tự xuất hiện: Định nghĩa 3.1, Định lý 3.2, Hệ quả 3.3, Bổ đề 3.4.

\clearpage
